\section{Commercial survival software and repeated banking interventions}
\label{exec15:commercial-survival}

A servicing team observes default, prepayment, hardship contact and changing
arrears over a loan's life. It asks whether repeated assistance will improve
customer outcomes and the value of the portfolio. Several familiar software
operations sound relevant: a survival curve, a Cox regression with changing
covariates, an average treatment effect on survival time, a competing-event
model and a latent growth model. They do not answer the same question.
An adequate comparison begins with the decision, eligible population, common
time origin, event definitions and intervention history. Only then can the
proposed library decide which commercial operation is a valid comparator.

This bounded review develops eight families from official technical text:
Stata, SAS/STAT, IBM SPSS, JMP, Minitab, NCSS, Statistica and Mplus. MATLAB
remains a ninth candidate because both attempted official method pages were
unavailable. Several sources identify an edition; the exact NCSS and Mplus
editions and the Minitab update remain unresolved. A mutable manual acquired
after 7 October 2026 does not establish its exact content at that literature
cutoff. Consequently the comparison supports its inspected method claims,
not an exhaustive historical market census or installed-product qualification.
No commercial engine was executed, and no price, deployment entitlement or
performance advantage is inferred from documentation
\citep{exec14_survival_survey}.

\subsection{Time, events and causal targets belong in the interface}

Let $T^g$ be the first terminal-event time under a declared feasible policy
$g$, and $J^g$ the event type, such as default or prepayment. A first-event
survival probability $P(T^g>t)$ differs from the default cumulative incidence
$F_D^g(t)=P(T^g\leq t,J^g=D)$. The restricted mean time without either event
is $E[\min(T^g,\tau)]$, for a specified horizon $\tau$. It is neither the
unrestricted mean default time nor a ratio of conditional hazards. A policy
may change prepayment as well as default, so changing the target can change
the apparent benefit. These objects are proposed contract fields, not
automatically identified quantities.

For a discrete two-period construction, write $d_k^g$ and $p_k^g$ for the
conditional probabilities of default and prepayment during period $k$ among
loans still active at its start. They are mutually exclusive, with
$d_k^g+p_k^g\leq1$. Set $S_0^g=1$. Direct conditioning gives
\[
 S_k^g=S_{k-1}^g(1-d_k^g-p_k^g),\qquad
 F_D^g(K)=\sum_{k=1}^K S_{k-1}^g d_k^g.
\]
The analogous sum with $p_k^g$ gives prepayment incidence. The recursion is
an accounting identity within the stipulated policy law. Identification of
that law from bank observations requires additional causal and observation
assumptions. In particular, replacing the recursion by
$1-\prod_k(1-d_k^g)$ removes prepayment while retaining the same default
hazards. That is a different construction; the unchanged-hazard premise
does not follow from observing the original competing-event process
\citep{exec14_survival_exact}.

Suppose both periods have default probability $1/10$ under two generated
policies. Policy A has prepayment probability $2/5$, while B has probability
$1/10$. After two periods A has default incidence $.15$, prepayment incidence
$.60$ and active probability $.25$. B has corresponding probabilities $.18$,
$.18$ and $.64$. Removing prepayment and retaining the default hazard gives
$.19$ under either policy. All three default numbers answer different
questions. A lower observed default incidence under A can arise because more
loans leave through prepayment, even though the conditional default hazard
is identical. Full enumeration of both period outcomes agrees with the
recursion \citep{exec14_survival_exact}.

Value need not rank the policies by default incidence. Stipulate USD~10 per
active loan at the start of each period, USD~100 lost at default and USD~20
lost at prepayment, with no discounting. Expected active periods are $1.5$
and $1.8$, so values are
\[
 V_A=10(1.5)-100(.15)-20(.60)=-12,\qquad
 V_B=10(1.8)-100(.18)-20(.18)=-3.6.
\]
The policy with fewer defaults has lower value under these exposed prices.
These are authored scenarios, not estimated bank effects or claims that
prepayment always causes a loss. The interface should return the joint event
accounting and the declared value map, allowing different accepted prices
without changing the scientific target \citep{exec14_survival_exact}.

\subsection{Stata offers causal survival comparators with explicit limits}

The inspected \texttt{stteffects} manuals document ATEs, effects on the treated
and potential-outcome mean survival times through regression adjustment,
inverse-probability weighting and combined methods. These are substantial
existing causal-estimation capabilities. Comparing the proposed library only
with an unadjusted survival plot would understate the commercial alternative.
The manuals also expressly exclude delayed entry, time-varying covariates and
multiple-record data from this command family. A documented baseline-treatment
survival estimator therefore cannot be silently relabeled as an estimator of
a repeated collections policy \citep{exec14_survival_stata}.

Three scientific conditions remain distinct: conditional independence of
treatment and potential outcomes given the supplied covariates, treatment
overlap and correct adjustment for censoring. The documentation distinguishes
likelihood-adjusted censoring in LAC-IPWRA from weighted-adjusted censoring
in WAC-IPWRA. In the inspected implementation, weighted censoring requires
random, correctly modeled censoring and a censoring process that does not
vary by treatment level. The choice is consequential for a servicing example
where treatment changes the chance that records disappear. Calling an
estimator doubly robust does not remove an identification premise or extend
the documented censoring regime \citep{exec14_survival_stata}.

The proposed adapter should retain the chosen estimand, treatment/control
labels, outcome distribution, treatment model and censoring adjustment as
separate fields. It should export propensity scores, censor-survival
probabilities, conditional means, covariance type and the actual estimation
sample. The manual documents an overlap tolerance that can stop estimation
and an option identifying offending observations. It also distinguishes
robust, clustered, bootstrap and jackknife uncertainty. A refusal at an
overlap boundary belongs in the result, not in an omitted request. Balance
and overlap diagnostics examine supplied observed covariates; they do not
authenticate an unmeasured-confounding assumption
\citep{exec14_survival_stata}.

Stata's separate \texttt{stcrreg} command models the subhazard of a selected
event with specified competing events. It permits single or multiple records
but not multiple failures per subject. Its event codes, time interactions,
cluster settings and cumulative-incidence interpretation require a different
conversion contract from \texttt{stteffects}. Recurrent missed payments and
terminal default are different event structures even when the raw table has
the same number of rows. The first proposed commercial trial should use a
supported fixed-treatment case and a separately specified competing-event
case; it should not claim that either demonstrates arbitrary dynamic-policy
evaluation \citep{exec14_survival_stata_competing}.

\subsection{SAS and JMP illustrate why competing-event outputs differ}

SAS/STAT 14.1 PHREG documents counting-process input in intervals $(T_1,T_2]$,
with covariates held fixed within each record and event status at the end.
Its examples cover repeated events and delayed entry. This is useful for
representing servicing records faithfully, provided subject identity, actual
entry time and covariate availability are preserved. The input representation
does not itself identify the effect of intervening on a history. It specifies
who belongs in a model's risk set \citep{exec14_survival_sas}.

The same edition documents a distinct subdistribution-hazard pathway through
\texttt{EVENTCODE}. Subjects with competing events remain in modified risk
sets with weights; they should not be exported as ordinary event-free subjects.
The manual specifies a sandwich covariance, an explicit cluster-ID route,
cumulative-incidence predictions and pointwise interval settings. It also
states that several options and statements are ignored in this mode, including
\texttt{ASSESS}, \texttt{BAYES} and \texttt{RANDOM}. A generic wrapper that
accepts every option and labels it effective would misdescribe the analysis.
The adapter must preserve the effective procedure configuration and distinguish
a pointwise band from a simultaneous or policy-value guarantee
\citep{exec14_survival_sas}.

JMP 18.1's inspected Competing Causes page describes cause-specific survival
estimation by censoring events from other causes and states independence for
its cause-wise interpretation. Its Omit Causes control recalculates the
display to illustrate removal of selected causes. In the banking construction,
removing prepayment and holding default hazards unchanged yields $.19$, while
actual default incidences are $.15$ and $.18$. Thus a cause-removal display
is not interchangeable with the observed cumulative default probability.
Nor does the control establish that a bank can remove prepayment without
changing default mechanisms. This is a scoped comparison of the documented
platform, not an assertion that every JMP survival capability shares that
interpretation \citep{exec14_survival_jmp,exec14_survival_exact}.

The proposal should therefore avoid importing a column merely named
``survival'' or ``risk.'' Its descriptor needs the event set, horizon,
conditional population, risk-set convention, cause-removal premise where
used, and uncertainty definition. SAS's subdistribution output and the
inspected JMP cause-wise display are useful alternatives for different
questions. A matched experiment first harmonizes the question and only then
compares numerical estimates, diagnostics and workflow cost.

\subsection{Time-dependent regression must preserve the decision clock}

IBM SPSS Statistics 30 documents the system time variable \texttt{T\_},
both for explicit time interactions and for segmented measured covariates.
Its weekly example combines logical interval indicators with measurements.
Statistica 14.2 separately offers survival-time inputs, explicit time-function
covariates and counting-process start/stop inputs. These are useful ways to
represent changing records, but a change of input form is not evidence of a
causal dynamic-regime estimator. Measurement time, availability time and the
action time must still be checked independently
\citep{exec14_survival_ibm,exec14_survival_statistica}.

Consider 100 stipulated loans at a common initial eligibility date. Twenty
default before day two. Among the 80 survivors, 40 receive assistance at day
two and 40 do not; ten in each group subsequently default. Labeling exposure
from the initial date by whether assistance ever occurs gives risks $10/40$
and $30/60$, an apparent difference of minus 25 percentage points. The
future-treated category necessarily excludes the twenty early defaults.
A day-two landmark comparison gives $10/40$ in each arm, but now concerns
the survivor population. Neither rewriting the label nor dropping early events
recovers a baseline assignment-policy effect. The construction diagnoses a
clock and denominator error; it is not an estimate of Cox-regression bias
\citep{exec14_survival_exact}.

A converter should split records at declared changes while preserving which
covariate value was known immediately before the decision. In a $(s,t]$
record, the end event and any new information at $t$ must follow an explicit
ordering convention. A measurement collected after an event cannot be carried
back into its risk interval. A policy scheduled for day two needs a definition
for loans that default earlier. These are scientific design choices, not
formatting repairs; the software should reject ambiguity rather than choose
an ordering from the row sort.

Minitab's inspected counting-process method table makes several conversion
checks concrete. It distinguishes fixed predictors from time-dependent
repeated measures, records categorical coding and tie adjustment, distinguishes
robust from model covariance, and describes stratification and proportional hazards diagnostics. It also documents excluded rows with missing or
nonpositive response times and an error when entry exceeds exit. Consequently
the input row count is insufficient: a proposed adapter should reconcile
subjects, intervals, events and exclusions, and preserve the effective
covariance and tie method. This is documentation-based design; exact edition
and runtime parity still need qualification \citep{exec14_survival_minitab}.

\subsection{Sequential confounding is not removed by a time column}

An exact two-decision construction shows the additional causal work. Let
$A_0$ be an early assistance action randomized with probability $1/2$.
Let $L$ be later arrears, with probabilities $.8$ under $A_0=0$ and $.2$
under $A_0=1$. The later action $A_1$ has probability $.8$ when $L=1$ and
$.2$ when $L=0$. Define final default probability by
\[
 P(Y=1\mid L=l,A_1=a)=.1+.4l-.05a.
\]
Independent randomizers generate the stated conditional probabilities. This
fully specified law supplies consistency, sequential exchangeability and
positive treatment support for this example. None is inferred from a fitted
hazard model or a software diagnostic \citep{exec14_survival_exact}.

Marginally $P(L=1)=1/2$. Later assistance selects higher-risk histories:
$P(L=1\mid A_1=1)=.8$ and $P(L=1\mid A_1=0)=.2$. Hence observed default
risks are $.37$ with $A_1=1$ and $.18$ with $A_1=0$, despite the negative
conditional action effect. Under fixed policies, direct summation instead gives
\[
 \begin{split}
 P(Y^{0,0}=1)&=.2(.1)+.8(.5)=.42,\\
 P(Y^{1,1}=1)&=.8(.05)+.2(.45)=.13.
 \end{split}
\]
Their difference is $-.29$, while the observed late-action difference is
$+.19$. The two policy values are also reproduced by averaging the default indicator $Y$ over the complete
sixteen-state law with weights
\[
 \frac{\mathbf 1\{A_0=a_0,A_1=a_1\}}
 {P(A_0=a_0)P(A_1=a_1\mid L)}.
\]
This is a population identity for the authored law, not a fitted weighting
study or an asserted finite-sample coverage result
\citep{exec14_survival_exact}.

Conditioning on $L$ can remove confounding of the later action while also
holding fixed a variable changed by the earlier action. In this construction,
the controlled $A_0$ difference at fixed $(L,A_1)$ is zero, although the
two-regime contrast is not. A single coefficient from a regression containing
$A_0,L,A_1$ cannot be interpreted as the total effect of the repeated policy
without further calculation and assumptions. The proposed library needs an
explicit longitudinal target, history-indexed assignment/observation laws
and a valid standardization or other justified procedure. A counting-process
adapter can supply faithful data to such work; it is not a substitute for it.

Mplus's inspected Chapter 6 offers distinct discrete-time survival, Cox,
parametric baseline-hazard and latent-factor survival examples. In the
discrete example, post-event periods and dropout can share a missing-value
code. A bank adapter should retain their different reasons before conversion.
The latent-factor example makes numerical integration and its computational
burden explicit. Its measurement model is a substantive alternative when
supported, but fitting a latent factor does not authenticate causal confounding
control. These examples justify a narrow model-and-output comparator, while
the chapter's exact edition, intended runtime and rights remain unresolved
\citep{exec14_survival_mplus}.

\subsection{Censoring and failed uncertainty calculations remain outcomes}

A small horizon-observation example separates an outcome model from its
observed cases. Half the population is low risk with default probability $.1$;
half is high risk with probability $.6$. All low-risk outcomes and one quarter
of high-risk outcomes are observed. Observation is stipulated independent
of default conditional on this recorded stratum. Full-cohort risk is
$.5(.1)+.5(.6)=.35$, while complete-case risk is
\[
 \frac{.5(.1)+.5(.25)(.6)}{.5+.5(.25)}=.20.
\]
Inverse observation weighting restores $.35$ under the known probabilities.
It cannot do so by arithmetic alone if the stratum is unavailable, the
probabilities are wrong or observation depends on unrecorded future outcomes.
This finite missingness construction motivates the analogous checks for a
time-indexed censoring model; it does not establish a continuous-time estimator
\citep{exec14_survival_exact}.

NCSS's Cox manual supplies a useful reporting comparator even without a
qualified integration. Its run summary contains processed and estimation
rows, likelihood iterations, achieved convergence and completion status.
The described stopping statistic is the maximum relative coefficient change
on the last iteration. A finite number in a coefficient table is therefore
not enough to classify a run as converged. The proposal should preserve
these diagnostics and distinguish numerical nonconvergence, a scientifically
unidentified target and an unsupported method. The source is an undated
official chapter; it supports the inspected reporting semantics, not a
claim about a particular installed NCSS release \citep{exec14_survival_ncss}.

Uncertainty calculations need their own failure denominator. Suppose eight
subjects include only one with a required rare treatment history. An ordinary
eight-draw subject bootstrap omits that subject with probability
$(7/8)^8\simeq.3436$. An unregularized mean for that history is undefined in
those resamples. Discarding them and silently drawing replacements conditions
the replication distribution on containing the history. It is not the
unconditional bootstrap originally declared. This exact occupancy calculation
does not claim a 34.36\% failure rate for any vendor command: it specifies a
stress case for the proposed wrapper. A cluster bootstrap, stratified design,
regularization or a restricted target is a substantive alternative with its
own assumptions, not a license to hide failures \citep{exec14_survival_exact}.

\subsection{Build a target-preserving comparator before a product claim}

The six-operation reference revision examined here handled declared graph
syntax, a supplied-DAG backdoor case, supplied-moment inversion, deterministic
sensitivity sets and declared observation clocks. It did not expose fitted
survival, competing-risk or dynamic-regime estimation. The later servicing
study in Chapter~\ref{gap29:servicing} supplies bounded fitted research;
Chapter~\ref{gap29:public} demonstrates fitted-summary inference across
interfaces. Neither result retroactively changes this probe or establishes a
general raw-history survival service. The inspected
commercial procedures already supply serious estimation, diagnostics and
reporting alternatives; the proposed differentiator must be demonstrated in
target preservation, supported causal scope and transparent uncertainty,
rather than assumed from having an agent interface
\citep{exec14_survival_survey}.

The immediate route is a public/generated comparator contract. Begin with a
fixed-treatment survival target supported by the chosen procedure; add a
competing-event case with separately declared incidence and cause-removal
targets. Verify subject identity, time units, entry/exit boundaries, event
codes, treatment levels, selected rows and output units. Then inspect complete
failure trajectories, censoring diagnostics and uncertainty settings. Repeated
policies enter a separate experiment with an explicit sequential identification
argument and exact reference law. No unsupported request should be routed to
the nearest command by name.

Stata and SAS warrant consideration as licensed numerical comparators;
IBM, Statistica and Minitab supply useful clock, representation and diagnostic
comparisons; JMP clarifies cause-removal semantics; NCSS exposes convergence
reporting; and Mplus supplies survival/measurement-model alternatives.
These are proposed uses of documented paths. Proprietary code should not be
ported under an assumed license. A future adapter or direct-use decision
requires the selected edition, authorized execution mode, permitted exports
and applicable redistribution and benchmark-publication terms. MATLAB and
uninspected extensions stay unresolved. Missing evidence is not a negative
feature finding \citep{exec14_survival_survey}.

Only after the public/generated prototype preserves these targets and makes
the intended comparison useful should a paid runtime study or authentic bank
experiment be commissioned. Its cost includes the license, conversion,
fitting, uncertainty replication, failed attempts, review and maintenance.
No current quotation or bank ROI is available from this review. That is a
specific later input, while the present proposal can already specify the
scientific comparisons, acceptance cases and unsupported uses.
\label{exec15:commercial-survival-end}
